Active context acquisition is often treated as inference from extra experience. In contact-rich material manipulation, however, the action that reveals hidden dynamics also changes the state that the downstream policy must control. We study this information--intervention tradeoff with \PTA{}, a deliberately minimal framework that executes a fixed probe, encodes the resulting proprioceptive and simulator-derived aggregate-particle trace into a latent task-conditioning vector, and conditions a residual policy on that vector. We evaluate on a Genesis edge-push diagnostic benchmark spanning sand, snow, and elastoplastic media; all methods train on sand only and transfer zero-shot. The empirical picture is asymmetric. On OOD elastoplastic, \PTA{} improves transfer efficiency by 14.7 percentage points (60.7\% vs.\ 46.0\%) and reduces spillage by 14.7 percentage points (39.3\% vs.\ 54.0\%) over reactive PPO; ablations suggest that both the fixed probe and latent conditioning contribute to this split-specific gain. On ID sand, two sand parameter shifts, and snow, the same probe-conditioned policy underperforms by 13--22 percentage points. We therefore do not claim broad cross-material robustness. Instead, the pattern is consistent with, but does not validate, a recoverable-response interpretation: learned context appears useful when probe-induced state shift relaxes on the task timescale, and harmful when the intervention leaves persistent disruption. \PTA{} is best read as a diagnostic tool for studying when active context acquisition is physically worth doing under hidden dynamics.
